\section{Discussion}

\subsection{A unified account of variability, sparseness, and efficiency}

The results presented above offer a unified perspective on a longstanding puzzle: how energy-efficient computation coexists with highly variable neuronal firing.  By formalising the energy-constrained coding problem on a Gibbs manifold, we show that sparseness and variability are intrinsically coupled descriptors of a single thermodynamic state, linked by the identity $S = \mathrm{CV}^2/(1+\mathrm{CV}^2)$ and jointly governed by the marginal efficiency $\beta$.

Optimal codes lie along a one-dimensional Pareto frontier in the $H$--$C$ plane \cite{laughlin_energy_2001, Harris2012, Shoval2012}.  When all three systems---V1 recordings, single-neuron plasticity, and recurrent network adaptation---are projected onto this plane (Fig.~6), their dynamics converge onto the same theoretical boundary, with ICE consistently exceeding $\beta$.  V1 spontaneous activity occupies the steepest segment, suggesting that the irregular resting state maintains maximum channel capacity within the prevailing energy budget \cite{Raichle2001}.

These findings complement recent evidence that energy constraints impose low-dimensional structure at the level of intrinsic neuronal parameters \cite{Sommer2026}, as well as classical information--rate analyses showing sublinear gains with firing rate \cite{Rieke1997,Sterling2015}.  Our framework addresses a different question---how a fixed budget constrains the entire \emph{distribution} of responses---and reveals that the sublinear information--rate relationship is a natural consequence of operating near the Pareto-optimal frontier.

\subsection{Thermodynamic interpretation and connections to the free energy principle}

The Gibbs form of the optimal distribution invites a thermodynamic reading.  The Lagrange multiplier $\beta$ is formally identical to an inverse temperature, and the manifold traces the states of a system in thermal equilibrium with a metabolic ``reservoir.''  This analogy is more than metaphorical: the variational structure underlying our framework---maximising entropy subject to an expected-cost constraint---is mathematically equivalent to the constrained free energy minimisation $F = C - \beta^{-1}H$, where $F$ is a neural analogue of Helmholtz free energy.  The condition $dF/dp = 0$ recovers the Gibbs distribution, and the manifold parameterised by $\beta$ maps directly onto an equilibrium thermodynamic state space.

This connection resonates with, but is distinct from, the free energy principle (FEP) of Friston \cite{Friston2006}.  The FEP minimises variational free energy to maintain homeostasis and model the external world; our framework minimises a thermodynamic free energy to maximise information throughput per unit cost.  Both select exponential-family distributions, but the constraints differ: the FEP constrains divergence from a generative model, whereas our framework constrains metabolic expenditure.  The two perspectives may converge in settings where metabolic efficiency and accurate world-modelling are jointly optimised---an intersection that merits further theoretical investigation.

\subsection{From biological mechanisms to engineering principles}

The geometry of the Gibbs manifold is fixed by the optimisation problem, independent of implementation.  Synaptic plasticity and E-I balance determine \emph{where} the system operates along it, not the manifold's shape.  Our models show that local plasticity rules---inhibitory STDP and Hebbian/anti-Hebbian learning---suffice to steer circuits toward the high-CV, high-efficiency regime \cite{43}, consistent with theories linking E-I balance to criticality \cite{Liang2025prl} but adding a thermodynamic dimension.

This biological principle translates into an engineering design rule.  Replacing hand-tuned sparsity priors with CV-maximisation recovers V1-like receptive fields with higher ICE and lower reconstruction error than classical sparse coding, demonstrating that targeting a macroscopic statistical signature provides a direct path to efficient representations.  For neuromorphic hardware operating under strict power budgets, this finding suggests that monitoring and maximising CV at the population level may be a more effective and hardware-agnostic objective than enforcing explicit sparsity penalties, which require distributional assumptions that may not hold in practice.

\subsection{Limitations and testable predictions}

Several limitations qualify our conclusions.  First, the theory assumes convex cost functions and a deterministic encoding approximation ($H(r|s)\approx 0$); when noise entropy is non-negligible, the Gibbs solution provides an upper bound on achievable efficiency rather than an exact optimum.  Second, the predicted monotonic ICE--CV relationship applies to \emph{population} CV, while single-neuron $\mathrm{CV}_L$ can saturate due to biophysical constraints such as the refractory period and membrane time constant, as demonstrated in our single-neuron simulations.  Third, the V1 data were recorded under sufentanil anesthesia with drifting grating stimuli; whether the same manifold structure holds under awake, free-viewing conditions with naturalistic stimuli remains to be verified.

The framework generates testable predictions: (i) acutely reducing the energy budget (e.g., cortical cooling or pharmacological suppression of oxidative metabolism) should decrease population CV and sparseness in tandem along the manifold, following the trajectory predicted by the $H$--$C$ curve; (ii) brain regions with tighter metabolic budgets (e.g., higher neuron density relative to vascular supply) should show higher CV and ICE; (iii) developmental sharpening of E-I balance should drive circuits toward the high-CV regime, a trajectory that may reverse with aging or neurodegeneration as inhibitory circuits deteriorate.  Analytical derivations supporting these predictions, including the monotonicity proofs for ICE--CV and collision probability--CV relationships, are provided in Supplementary Sections~1--2.

\subsection{Conclusion}

Metabolic constraints are not passive boundaries but active forces shaping neural codes.  Under energy-constrained optimisation, high variability emerges as a mathematical inevitability---the regime where information per unit cost is maximised, sparseness is intrinsically guaranteed, and robustness is highest.  CV serves as a single coordinate organising these properties along the Gibbs manifold, reframing cortical ``noise'' as the signature of thermodynamic optimality.  This perspective bridges neuroscience and engineering: in biological circuits, E-I balance and synaptic plasticity implement the navigation toward optimal CV; in artificial systems, directly targeting high CV offers a hardware-agnostic design principle for energy-efficient computation.
